\textbf{Phần 1. Trắc nghiệm lựa chọn}
\setcounterpageref{ex}{0}

\begin{ex} %[4D12V65]
Muốn tìm phân số $\dfrac{a}{b}$ của số $M$, ta tính
\choice
{$M : a \times b$}
{$M + \dfrac{a}{b}$}
{\True $M \times \dfrac{a}{b}$ (hoặc lấy $M : b \times a$)}
{$M - \dfrac{a}{b}$}
\loigiai{
Ta lấy số $M$ nhân với phân số $\dfrac{a}{b}$.\\
\textbf{Đáp án của câu hỏi là C}
}
\end{ex}

\begin{ex} %[4D12V65]
$\dfrac{2}{3}$ của $18$ là
\choice
{\True $12$}
{$6$}
{$9$}
{$15$}
\loigiai{
$18 \times \dfrac{2}{3} = 12$.\\
\textbf{Đáp án của câu hỏi là A}
}
\end{ex}

\begin{ex} %[4D12V65]
$\dfrac{3}{5}$ của $100\text{ kg}$ là
\choice
{$50\text{ kg}$}
{$70\text{ kg}$}
{$40\text{ kg}$}
{\True $60\text{ kg}$}
\loigiai{
$100 \times \dfrac{3}{5} = 60\text{ kg}$.\\
\textbf{Đáp án của câu hỏi là D}
}
\end{ex}

\begin{ex} %[4D12V65]
$\dfrac{3}{4}$ của $1$ giờ là bao nhiêu phút?
\choice
{$40$ phút}
{\True $45$ phút}
{$50$ phút}
{$30$ phút}
\loigiai{
$60 \times \dfrac{3}{4} = 45$ phút.\\
\textbf{Đáp án của câu hỏi là B}
}
\end{ex}

\begin{ex} %[4D12V65]
Một lớp học có 36 học sinh, trong đó $\dfrac{4}{9}$ số học sinh là nữ. Số học sinh nữ của lớp đó là
\choice
{$14$ học sinh}
{$18$ học sinh}
{\True $16$ học sinh}
{$20$ học sinh}
\loigiai{
$36 \times \dfrac{4}{9} = 16$ học sinh.\\
\textbf{Đáp án của câu hỏi là C}
}
\end{ex}

\begin{ex} %[4D12V65]
$\dfrac{5}{6}$ của $240\text{ m}$ là
\choice
{\True $200\text{ m}$}
{$180\text{ m}$}
{$220\text{ m}$}
{$160\text{ m}$}
\loigiai{
$240 \times \dfrac{5}{6} = 200\text{ m}$.\\
\textbf{Đáp án của câu hỏi là A}
}
\end{ex}

\begin{ex} %[4D12V65]
Một cửa hàng có $150\text{ kg}$ gạo, đã bán được $\dfrac{2}{5}$ số gạo đó. Số gạo cửa hàng còn lại là
\choice
{$60\text{ kg}$}
{$100\text{ kg}$}
{$80\text{ kg}$}
{\True $90\text{ kg}$}
\loigiai{
Đã bán: $150 \times \dfrac{2}{5} = 60\text{ kg}$.\\
Còn lại: $150 - 60 = 90\text{ kg}$.\\
\textbf{Đáp án của câu hỏi là D}
}
\end{ex}

\begin{ex} %[4D12V65]
Một đàn gà có 48 con. Số gà trống bằng $\dfrac{3}{8}$ tổng số gà. Số gà mái là
\choice
{$18$ con}
{\True $30$ con}
{$24$ con}
{$28$ con}
\loigiai{
Gà trống: $48 \times \dfrac{3}{8} = 18$ con.\\
Gà mái: $48 - 18 = 30$ con.\\
\textbf{Đáp án của câu hỏi là B}
}
\end{ex}

\begin{ex} %[4D12V65]
Một mảnh đất hình chữ nhật có chiều dài $40\text{ m}$, chiều rộng bằng $\dfrac{3}{4}$ chiều dài. Diện tích mảnh đất đó là
\choice
{$1000\text{ m}^2$}
{$1400\text{ m}^2$}
{\True $1200\text{ m}^2$}
{$1600\text{ m}^2$}
\loigiai{
Chiều rộng: $40 \times \dfrac{3}{4} = 30\text{ m}$.\\
Diện tích: $40 \times 30 = 1200\text{ m}^2$.\\
\textbf{Đáp án của câu hỏi là C}
}
\end{ex}

\begin{ex} %[4D12V65]
Bạn Nam có $120\,000$ đồng. Nam mua sách hết $\dfrac{2}{3}$ số tiền đó. Số tiền còn lại của Nam là
\choice
{\True $40\,000$ đồng}
{$30\,000$ đồng}
{$50\,000$ đồng}
{$80\,000$ đồng}
\loigiai{
Tiền mua sách: $120\,000 \times \dfrac{2}{3} = 80\,000$ đồng.\\
Còn lại: $120\,000 - 80\,000 = 40\,000$ đồng.\\
\textbf{Đáp án của câu hỏi là A}
}
\end{ex}

\begin{ex} %[4D12V65]
Một quãng đường dài $120\text{ km}$. Ô tô đã đi được $\dfrac{5}{8}$ quãng đường. Quãng đường ô tô còn phải đi tiếp là
\choice
{$40\text{ km}$}
{$50\text{ km}$}
{$75\text{ km}$}
{\True $45\text{ km}$}
\loigiai{
Đã đi: $120 \times \dfrac{5}{8} = 75\text{ km}$. Còn lại: $120 - 75 = 45\text{ km}$.\\
\textbf{Đáp án của câu hỏi là D}
}
\end{ex}

\begin{ex} %[4D12V65]
Một rổ cam có 60 quả. Lần thứ nhất bán $\dfrac{1}{3}$ số cam, lần thứ hai bán $\dfrac{1}{2}$ số cam CÒN LẠI. Số cam còn lại trong rổ sau 2 lần bán là
\choice
{$15$ quả}
{\True $20$ quả}
{$25$ quả}
{$10$ quả}
\loigiai{
Lần 1 bán: $60 \times \dfrac{1}{3} = 20$ quả $\Rightarrow$ còn 40 quả.\\
Lần 2 bán: $40 \times \dfrac{1}{2} = 20$ quả.\\
Sau 2 lần còn: $40 - 20 = 20$ quả.\\
\textbf{Đáp án của câu hỏi là B}
}
\end{ex}

\textbf{Phần 2. Trắc nghiệm đúng sai}
\setcounterpageref{ex}{0}

\begin{ex} %[4D12V65]
Xét tính đúng sai về tìm phân số của một số:
\choiceTF
{\True $\dfrac{3}{4}$ của 24 là 18.}
{\True $\dfrac{2}{5}$ của $1\text{ tấn}$ là $400\text{ kg}$.}
{$\dfrac{1}{4}$ thế kỉ là $20$ năm.}
{\True $\dfrac{1}{4}$ thế kỉ là $25$ năm.}
\loigiai{
\begin{itemize}
    \item $24 \times 3/4 = 18$. Đúng.
    \item $1000 \times 2/5 = 400\text{ kg}$. Đúng.
    \item $100 : 4 = 25$ năm $\neq 20$ năm. Ý 3 Sai, Ý 4 Đúng.
\end{itemize}
\begin{itemize}
    \item a) \textbf{Đúng}.
    \item b) \textbf{Đúng}.
    \item c) \textbf{Sai}.
    \item d) \textbf{Đúng}.
\end{itemize}
}
\end{ex}

\begin{ex} %[4D12V65]
Một khối lớp 4 có 160 học sinh, trong đó có $\dfrac{3}{5}$ số học sinh đạt học lực Xuất sắc, $\dfrac{1}{4}$ đạt học lực Tốt, còn lại là học lực Hoàn thành. Xét tính đúng sai:
\choiceTF
{\True Số học sinh Xuất sắc là 96 em.}
{\True Số học sinh Tốt là 40 em.}
{Số học sinh Hoàn thành là 30 em.}
{\True Số học sinh Hoàn thành là 24 em.}
\loigiai{
\begin{itemize}
    \item Xuất sắc: $160 \times 3/5 = 96$ em. Đúng.
    \item Tốt: $160 \times 1/4 = 40$ em. Đúng.
    \item Hoàn thành: $160 - (96 + 40) = 24$ em $\neq 30$ em. Ý 3 Sai, Ý 4 Đúng.
\end{itemize}
\begin{itemize}
    \item a) \textbf{Đúng}.
    \item b) \textbf{Đúng}.
    \item c) \textbf{Sai}.
    \item d) \textbf{Đúng}.
\end{itemize}
}
\end{ex}

\textbf{Phần 3. Trả lời ngắn}
\setcounterpageref{ex}{0}

\begin{ex} %[4D12V65]
Tìm $\dfrac{3}{5}$ của $45$.
\shortans{27}
\loigiai{
\textbf{Đáp số:} $27$\\
$45 \times 3/5 = 27$.
}
\end{ex}

\begin{ex} %[4D12V65]
Tìm $\dfrac{5}{8}$ của $64\text{ kg}$.
\shortans{40}
\loigiai{
\textbf{Đáp số:} $40$\\
$64 \times 5/8 = 40\text{ kg}$.
}
\end{ex}

\begin{ex} %[4D12V65]
Tìm $\dfrac{2}{3}$ của $1$ giờ theo phút.
\shortans{40}
\loigiai{
\textbf{Đáp số:} $40$\\
$60 \times 2/3 = 40$ phút.
}
\end{ex}

\begin{ex} %[4D12V65]
Một mảnh đất hình chữ nhật dài $50\text{ m}$, chiều rộng bằng $\dfrac{3}{5}$ chiều dài. Tính diện tích mảnh đất theo mét vuông.
\shortans{1500}
\loigiai{
\textbf{Đáp số:} $1500$\\
Rộng: $30\text{ m}$. Diện tích: $50 \times 30 = 1500\text{ m}^2$.
}
\end{ex}

\begin{ex} %[4D12V65]
Một rổ có 40 quả táo, mẹ đã biếu bà $\dfrac{3}{8}$ số táo đó. Hỏi trong rổ còn lại bao nhiêu quả táo?
\shortans{25}
\loigiai{
\textbf{Đáp số:} $25$\\
Biếu bà: $40 \times 3/8 = 15$ quả. Còn lại: $40 - 15 = 25$ quả.
}
\end{ex}

\begin{ex} %[4D12V65]
Một bình nước có 24 lít. Người ta lấy ra lần 1 là $\dfrac{1}{4}$ lượng nước trong bình, lần 2 lấy tiếp $\dfrac{1}{3}$ lượng nước ban đầu. Tính lượng nước còn lại trong bình theo lít.
\shortans{10}
\loigiai{
\textbf{Đáp số:} $10$\\
Lấy ra: $24 \times (1/4 + 1/3) = 24 \times 7/12 = 14$ lít. Còn lại: $24 - 14 = 10$ lít.
}
\end{ex}
